$\phi_{\mathrm a}$ characterizes the fraction of the sliding side chain within the A component.
It is a fixed structural parameter set by molecular synthesis, and serves as the primary control variable for tuning the sliding effect in L--S TT copolymers.
Although L--S TT copolymers with different $\phi_{\mathrm a}$ all reduce to the same class of linear AB diblock copolymers in the limit of $f=0$, their conformational response depends on $\phi_{\mathrm a}$ once sliding is allowed.
The effect of sliding on this response is expected to be strongest at intermediate $\phi_{\mathrm a}$.
When $\phi_{\mathrm a}$ is small, the system corresponds to a short side chain sliding along a relatively long main-chain track; although the accessible sliding range is large, the short side chain has only a limited ability to regulate the overall conformational statistics.
When $\phi_{\mathrm a}$ is large, the sliding side chain is longer but is confined to a shorter main-chain track, so that the topological degree of freedom cannot be fully released.
Therefore, $\phi_{\mathrm a}$ effectively measures the ability of the system to use the sliding mechanism to regulate conformational entropy and interfacial structure.
A systematic study of the self-assembly behavior at different $\phi_{\mathrm a}$ therefore helps clarify how this conformational response influences microphase separation.


Compared with conventional copolymer systems, another important feature of L--S TT copolymers is the spontaneous response of the sliding degree of freedom $f$.
During self-assembly, $f$ is not externally prescribed; instead, the relative contributions of chains with different $f$ to the local A-segment density are determined self-consistently by the equilibrium fields, reflecting geometric constraints, interfacial curvature, chain stretching, and density-continuity requirements.
This local topological selection generates a spatially resolved distribution of topological preference, $P_{\mathcal R(\bm r)}(f)$.
As a result, sliding topology not only modifies the phase diagram and spinodal line, but also gives rise to spatially heterogeneous topology responses to interfacial curvature and local packing constraints.
Therefore, analyzing the spatial variation of $P_{\mathcal R(\bm r)}(f)$ and its relation to the local environment is essential for understanding not only how sliding regulates microphase formation through topological redistribution, but also how local topology selection responds to heterogeneous packing environments within different microphases.


In this physical picture, $\phi_{\mathrm a}$ sets the molecular capacity for topology regulation by controlling the relative side-chain length and sliding range, whereas $f$ remains a chain-level internal variable whose conditional distribution among locally sampled A segments depends on the surrounding environment.
Based on these two key variables, we organize the following results from global to local behavior and from thermodynamic properties to structural mechanisms, thereby revealing how sliding topology regulates microphase separation.
We first focus on the ODT at $\bar{\phi}_{\mathrm A}=0.50$.
The close agreement between the SCFT-predicted ODT and the RPA spinodal indicates that the onset of microphase separation in L--S TT copolymers is governed primarily by the quadratic instability of the homogeneous phase.

The entropic consequences of sliding are characterized using two complementary measures.
The effective conformational-entropy scale $\Delta S_{\mathrm{conf}}$, obtained from the RPA spinodal difference, characterizes the conformational response associated with the homogeneous instability, whereas the sliding-distribution entropy loss $\Delta S_{\mathrm{Slide}}$, evaluated from the local topology distributions, measures the overall departure of the sliding ensemble from a uniform distribution after ordering.

The resulting SCFT phase diagrams for $\phi_{\mathrm a}=0.25$, $0.50$, and $0.75$ establish how molecular architecture reshapes phase stability, while the Gyroid phase provides a representative example of how variations in $\phi_{\mathrm a}$ redistribute the spatial roles of the constituent blocks and shift the global preferred topology $f_{\Lambda}^{\ast}$.
In the O$^{70}$ phase, the evolution of $P_{\Omega}(f)$ across isodensity layers, together with the relation between local preferred topology and mean curvature, reveals a geometry-conditioned topology response that differs between A-core and A-shell regions.
The spherical BCC and FCC phases further show that pronounced spatial variations in preferred topology can coexist with broad and strongly overlapping local topology distributions.
This coexistence gives rise to topology homogenization, through which local topology regulation remains spatially responsive without requiring strong localization of the full sliding ensemble, thereby providing a route to accommodate heterogeneous filling environments associated with packing frustration.

The dependence on segregation strength is finally examined through the $N\chi$-dependent phase widths and the spatial evolution of the preferred topology.
The contrasting responses of BCC and FCC indicate that stabilization by sliding is not determined solely by interfacial complexity, but also by whether topology regulation remains active over a sufficiently broad spatial region to accommodate local packing and interfacial requirements.

Taken together, these results show that the conformational response near the homogeneous instability, the redistribution of the sliding ensemble after ordering, and the spatial regulation of topology within ordered microphases are coupled consequences of the same annealed sliding degree of freedom.


\PDFsubsection{Topology Selection and Entropic Effects near the ODT}
    \begin{table}[htbp]
	\centering
	\small
	\renewcommand{\arraystretch}{1.37}
	\caption{
		Definitions of the theoretical systems considered for L--S TT copolymers.
	}
	\begin{tabular}{
		>{\centering\arraybackslash}m{2.3cm}
		>{\raggedright\arraybackslash}m{3.0cm}
		>{\raggedright\arraybackslash}m{5.0cm}
		>{\centering\arraybackslash}m{2.0cm}
		>{\centering\arraybackslash}m{2.0cm}
	}
		\toprule
		System & Topology treatment & Physical meaning & Transition parameter & Selected topology \\
		\midrule
        
		RPA--Slide & Freely sliding topology & Spinodal instability of the actual sliding system & $N\chi_{\mathrm{sp}}^{\mathrm{Slide}}$ & -- \\
        
		RPA--MUT & Single topology channel $P_{\mathrm{RPA}}(f)=\delta(f-f_0)$ & First unstable fixed-topology mode & $N\chi_{\mathrm{sp}}^{\mathrm{MUT}}$ & $f_{\mathrm{RPA\text{-}MUT}}^\ast$ \\
        
		SCFT--Slide & Freely sliding topology & Effect of sliding freedom on the actual phase behavior & $N\chi_{\mathrm{ODT}}^{\mathrm{Slide}}$ & $f_{\Lambda}^{\ast}$ \\
        
		SCFT--Fix & Fixed topology $f=f_{\mathrm{RPA\text{-}MUT}}^\ast$ & SCFT ODT used to assess the RPA spinodal prediction for the most unstable topology & $N\chi_{\mathrm{ODT}}^{\mathrm{Fix}}$ & -- \\
        
		Blend & Equal-weight mixture of all fixed topologies $f$ & Reference system for fixed-topology mixing & -- & -- \\

		\bottomrule
	\end{tabular}
	\label{tab:TheoryModels}
\end{table}



Establishing the phase behavior of L--S TT copolymers requires identifying both the stability limit of the homogeneous phase and the onset of ordered microphase separation, characterized by the RPA spinodal and the SCFT ODT, respectively.
Their comparison distinguishes the weak-segregation instability of the homogeneous state from nonlinear effects that shift the transition in the ordered phase.\cite{Matsen1996Unifying,Sides2004Continuous,Wang2008Microphase,Li2020Largely,Li2021Simplicity}
We also examine the global preferred topology $f_\Lambda^\ast$ near the ODT, because the topology selected at the onset of instability provides a reference for the molecular architecture that initiates ordering.
To separate the effects of freely sliding and fixed-topology treatments, we compare the five theoretical systems summarized in Table~\ref{tab:TheoryModels}.
RPA--Slide and SCFT--Slide describe the freely sliding L--S TT system at the spinodal and ODT levels, whereas RPA--MUT and SCFT--Fix provide the corresponding fixed-topology references.

The Blend system is included only as an equal-weight mixture of fixed topologies, thereby distinguishing prescribed statistical mixing from the self-consistently selected topology of the freely sliding L--S TT system.


\begin{figure}[!t]
\centering
\includegraphics[width=\linewidth]{./Figure/ODT.png}
\caption{
    (a) Global topology distribution $P_{\Lambda}(f)$ in the lamellar phase at fixed $\bar{\phi}_{\mathrm A}=0.50$ and $\phi_{\mathrm a}=0.50$ for different $N\chi$ near the ODT.
    Symbols denote SCFT data, solid lines denote spline fits, and stars mark the preferred topology $f_{\Lambda}^{\ast}$.
    (b) $P_{\Lambda}(f)$ for different $\phi_{\mathrm a}$ at fixed $\bar{\phi}_{\mathrm A}=0.50$ and $N\chi=N\chi_{\mathrm{ODT}}^{\mathrm{Slide}}+0.5$.
    (c) Spinodal and ODT interaction parameters as functions of $\phi_{\mathrm a}$.
    The blue solid and red dashed lines show $N\chi_{\mathrm{sp}}^{\mathrm{MUT}}$ and $N\chi_{\mathrm{sp}}^{\mathrm{Slide}}$, while the green pluses and orange crosses show $N\chi_{\mathrm{ODT}}^{\mathrm{Fix}}$ and $N\chi_{\mathrm{ODT}}^{\mathrm{Slide}}$, respectively.
    The gray dashed line marks the linear-diblock reference $N\chi=10.495$.
    (d) Preferred topology as a function of $\phi_{\mathrm a}$, with $f_{\mathrm{RPA\text{-}MUT}}^\ast$ from RPA--MUT shown by the blue line and $f_{\Lambda}^\ast$ from SCFT--Slide by the orange crosses.
    (e) Effective conformational entropy scale $\Delta S_{\mathrm{conf}}/k_{\mathrm B}$, defined by the RPA spinodal difference $N\chi_{\mathrm{sp}}^{\mathrm{Slide}}-N\chi_{\mathrm{sp}}^{\mathrm{MUT}}$.
    The SCFT ODT difference $N\chi_{\mathrm{ODT}}^{\mathrm{Slide}}-N\chi_{\mathrm{ODT}}^{\mathrm{Fix}}$ is shown for comparison.
    (f) Sliding distribution entropy loss $\Delta S_{\mathrm{Slide}}/k_{\mathrm B}$.
}
\label{fig:ODT}
\end{figure}


Fig.~\ref{fig:ODT}a shows the evolution of the global topology distribution $P_\Lambda(f)$ with increasing $N\chi$ at fixed $\bar{\phi}_{\mathrm A}=0.50$ and $\phi_{\mathrm a}=0.50$.
At the ODT, $P_\Lambda(f)=1$, consistent with $P_{\mathrm{RPA}}(f)=1$ for RPA--Slide at the homogeneous saddle point; consequently, the freely sliding system shares the same spinodal instability as the equal-weight Blend.
Beyond the ODT, $P_\Lambda(f)$ develops a maximum at the preferred topology $f_\Lambda^\ast$, showing that topology selection emerges with microphase separation.
The uniform weighting in RPA--Slide does not identify a topology that preferentially contributes to the emerging ordered structure.
RPA--MUT provides this complementary reference by searching the fixed-topology fluctuation channels for the first unstable topology, $f_{\mathrm{RPA\text{-}MUT}}^\ast$, which is expected to dominate the initial ordering near the ODT.
Fig.~\ref{fig:ODT}b extends this comparison across $\phi_{\mathrm a}$ by showing $P_\Lambda(f)$ at $\bar{\phi}_{\mathrm A}=0.50$ and $N\chi=N\chi_{\mathrm{ODT}}^{\mathrm{Slide}}+0.5$, revealing how the preferred topology shifts systematically with molecular architecture.
The resulting $f_\Lambda^\ast$ is compared with $f_{\mathrm{RPA\text{-}MUT}}^\ast$ in Fig.~\ref{fig:ODT}d.
The corresponding global preferred and mean topology statistics are discussed in Supporting Information Sec.~\ntxref{SI-sec:TopologyStatistics}, particularly Sec.~\ntxref{SI-subsec:TopologyStatisticsGlobal}.
Their close agreement, including the star-to-linear topology-locking crossover at $\phi_{\mathrm a}\approx0.68$, supports RPA--MUT as a predictor of the topology selected at the onset of ordering~\cite{Zhao2024Phase}. 
Here, topology locking refers to the selection of a preferred topology through the position of the distribution maximum, which identifies the topology favored by the ordered state while the full sliding ensemble can remain broadly distributed.



As shown in Fig.~\ref{fig:ODT}c, $N\chi_{\mathrm{ODT}}^{\mathrm{Slide}}$ closely follows $N\chi_{\mathrm{sp}}^{\mathrm{Slide}}$, while $N\chi_{\mathrm{ODT}}^{\mathrm{Fix}}$ similarly follows $N\chi_{\mathrm{sp}}^{\mathrm{MUT}}$; the endpoint values at $\phi_{\mathrm a}=0$ and $1$ are evaluated directly from the equivalent linear AB-diblock model. 
The numerical determination of the ODT is described in Supporting Information Sec.~\ntxref{SI-sec:ODT}.
 These parallel trends indicate that, near the ODT, the onset of microphase separation is governed primarily by the instability of the homogeneous phase against composition fluctuations, rather than by nonlinear effects specific to the ordered state. 

Within the same RPA framework, RPA--Slide and RPA--MUT differ only in their treatment of $P_{\mathrm{RPA}}(f)$: the former retains the full sliding ensemble, whereas the latter selects the most unstable fixed-topology reference, resulting in different single-chain restoring responses.
The resulting spinodal difference provides an effective conformational entropy scale that reflects the strength of this restoring response near the homogeneous instability, as detailed in Supporting Information Sec.~\ntxref{SI-sec:Entropy}.

Accordingly,
\begin{equation}
    \frac{\Delta S_{\mathrm{conf}}}{k_{\mathrm B}}
    \equiv
    N\chi_{\mathrm{sp}}^{\mathrm{Slide}}
    -
    N\chi_{\mathrm{sp}}^{\mathrm{MUT}}.
\end{equation}
Fig.~\ref{fig:ODT}e shows $\Delta S_{\mathrm{conf}}$ obtained from the RPA spinodal difference.
Because the SCFT ODT closely follows the corresponding RPA spinodal, the difference $N\chi_{\mathrm{ODT}}^{\mathrm{Slide}}-N\chi_{\mathrm{ODT}}^{\mathrm{Fix}}$ provides a good approximation to $\Delta S_{\mathrm{conf}}$ near the transition.
The maximum in $\Delta S_{\mathrm{conf}}$ occurs near the crossover between the star-locked and linear-locked regimes, where the preferred topology changes most rapidly with $\phi_{\mathrm a}$.
This behavior is consistent with the physical picture described above: competition between conformations with long and short side chains is strongest at the locking transition, maximizing the conformational benefit of sliding in this intermediate regime.


% The effective conformational-entropy measure $\Delta S_{\mathrm{conf}}$ compares the limiting fluctuation treatments of RPA--MUT and RPA--Slide and thus characterizes the stabilization of the full sliding ensemble relative to the most unstable fixed-topology channel.
% Once microphase separation occurs, the sliding ensemble becomes topologically and spatially biased; its entropy loss must therefore be evaluated from the local topology distribution $P_{\mathcal R(\bm r)}(f)$.
% We define this sliding-entropy loss as
% \begin{equation}
%     \frac{\Delta S_{\mathrm{Slide}}}{k_{\mathrm B}}
%     =
%     \frac{1}{\bar{\phi}_{\mathrm A}}
%     \int_\Lambda
%     \frac{\mathrm d\bm r}{V}\,
%     \phi_{\mathrm A}(\bm r)
%     \int_0^1
%     \mathrm df\,
%     P_{\mathcal R(\bm r)}(f)
%     \ln P_{\mathcal R(\bm r)}(f).
% \end{equation}
% Here, $\phi_{\mathrm A}(\bm r)$ weights the local entropy loss by the A-segment population sampled by $P_{\mathcal R(\bm r)}(f)$, and $\bar{\phi}_{\mathrm A}=V^{-1}\int_\Lambda\phi_{\mathrm A}(\bm r)\,\mathrm d\bm r$ provides the normalization.
% Marginalizing over $\bm r$ yields the global distribution $P_\Lambda(f)$ and the corresponding $\Delta S_\Lambda/k_{\mathrm B}=\int_0^1P_\Lambda(f)\ln P_\Lambda(f)\,\mathrm df$, which retains the global topology bias but omits the entropy contribution from spatially heterogeneous topology selection.
% The difference, $I_{f;\mathcal R}=(\Delta S_{\mathrm{Slide}}-\Delta S_\Lambda)/k_{\mathrm B}$, measures this missing contribution.
% Thus, $\Delta S_{\mathrm{Slide}}$ is a distribution-based entropy loss of the ordered sliding state, whereas $\Delta S_{\mathrm{conf}}$ is inferred from the shift between the fixed-topology and fully sliding instability limits and serves as an effective conformational-stabilization scale.
% As shown in Fig.~\ref{fig:ODT}e,f, $\Delta S_{\mathrm{Slide}}$ remains substantially smaller than $\Delta S_{\mathrm{conf}}$, consistent with the broad topology distribution that persists after the onset of microphase separation.
% Beyond their different magnitudes, the three entropy measures also exhibit distinct maxima as functions of $\phi_{\mathrm a}$, revealing different levels of topological response.
% $\Delta S_{\mathrm{conf}}$ is governed by the competition between topology channels at the homogeneous instability, whereas $\Delta S_{\mathrm{Slide}}$ captures the strongest spatially resolved topology selection in the ordered phase.
% Spatial marginalization further shifts the maximum of $\Delta S_\Lambda$, because locally preferred topologies can partially cancel when integrated over $\bm r$; the strongest local topology selection therefore need not coincide with the strongest global topology bias.

The effective conformational-entropy scale $\Delta S_{\mathrm{conf}}$ characterizes the increased resistance of the full sliding ensemble to composition fluctuations relative to the most unstable fixed-topology reference.
After microphase separation, local topology selection changes the topology weights from the uniform reference distribution.
We quantify the local departure from uniform sliding by the Kullback--Leibler (KL) divergence
\begin{equation}
    D_{\mathcal R}(\bm r)
    =
    \int_0^1
    \mathrm df\,
    P_{\mathcal R(\bm r)}(f)
    \ln P_{\mathcal R(\bm r)}(f).
    \label{eq:LocalDivergence}
\end{equation}
Here, $D_{\mathcal R}=0$ corresponds to a uniform local topology distribution, whereas larger $D_{\mathcal R}$ indicates a stronger local departure from uniform sliding.
The corresponding sliding-distribution entropy loss is defined as
\begin{equation}
    \frac{\Delta S_{\mathrm{Slide}}}{k_{\mathrm B}}
    =
    \frac{1}{V\bar{\phi}_{\mathrm A}}
    \int_\Lambda 
    \phi_{\mathrm A}(\bm r) D_{\mathcal R}(\bm r) \, \mathrm d\bm r.
    \label{eq:SlideEntropy}
\end{equation}
Here, $\phi_{\mathrm A}(\bm r)$ weights the local entropy loss by the A-segment population.
The normalization is provided by $\bar{\phi}_{\mathrm A}=V^{-1}\int_\Lambda \phi_{\mathrm A}(\bm r)\,\mathrm d\bm r$.



The local distribution is retained because different spatial regions can favor different topologies, whose preferences may be obscured in the global distribution $P_\Lambda(f)$.
The statistical construction and the comparison with the global description are given in Supporting Information Sec.~\ntxref{SI-sec:Entropy}.
Fig.~\ref{fig:ODT}e,f therefore characterizes two complementary aspects of the sliding response: the effective conformational entropy scale near the homogeneous instability and the sliding distribution entropy loss in the ordered phase.
As shown in Fig.~\ref{fig:ODT}f, $\Delta S_{\mathrm{Slide}}$ peaks near $\phi_{\mathrm a}=0.5$, indicating the strongest departure of the local topology distribution from uniform sliding at the intermediate architecture.
Its maximum occurs at a different $\phi_{\mathrm a}$ value from that of $\Delta S_{\mathrm{conf}}$, showing that the two quantities probe distinct aspects of the same sliding degree of freedom.
Their responses remain coupled because topology redistribution and chain conformations are determined self-consistently by the same molecular fields, as discussed further in Supporting Information Sec.~\ntxref{SI-sec:Entropy}.


\PDFsubsection{SCFT Phase Behavior and Topological Adaptation}
    \begin{figure*}[!ht]
\centering
\includegraphics[width=\linewidth]{./Figure/Phase.png}
\caption{
    (a--c) Phase diagrams of L--S TT copolymers in the $N\chi$--$\bar{\phi}_{\mathrm A}$ plane for $\phi_{\mathrm a}=0.25$, $0.50$, and $0.75$, respectively.
    Red lines denote SCFT results for the L--S TT system, gray lines denote the reference phase diagram of linear AB diblock copolymers, and blue dashed lines denote the RPA spinodal boundaries~\cite{Zhao2024Phase}.
    Insets show representative three-dimensional component-density isosurfaces of the main ordered phases, including FCC, BCC, C, G, $\mathrm{O}^{70}$, and L.
}
\label{fig:PhaseDiagram}
\end{figure*}



Sliding reshapes the phase boundaries of L--S TT copolymers.

Fig.~\ref{fig:PhaseDiagram} presents the SCFT phase diagrams of L--S TT copolymers for representative sliding side-chain fractions, $\phi_{\mathrm a}=0.25,\,0.50,\,0.75$, together with the spinodal boundaries (blue dashed lines) and SCFT ODT boundaries, and compares them with the phase diagram of conventional linear AB diblock copolymers (gray lines).
Because L--S TT molecules undergo a transition from star topology locking to linear topology locking as $\phi_{\mathrm a}$ changes, the spinodal lines at $\phi_{\mathrm a}=0.25$ and $\phi_{\mathrm a}=0.50$ lie slightly below those of linear AB diblock copolymers, whereas at $\phi_{\mathrm a}=0.75$ they almost coincide.

The sliding ensemble also modifies the conformational resistance of the homogeneous phase to composition fluctuations, as reflected by $\Delta S_{\mathrm{conf}}$ in the preceding analysis. 
As a result, for all values of $\phi_{\mathrm a}$, the SCFT ODT boundaries of L--S TT copolymers lie significantly above the corresponding SCFT ODT boundaries of linear AB diblock copolymers and also above the RPA spinodal lines of the L--S TT system.


% 图~\ref{fig:PhaseDiagram} 的另一重要特征是，当 $\phi_\mathrm{a}\neq 0$ 时，相图结构呈现出显著的非对称性，各相区不再关于 $\bar{\phi}_\mathrm{A}=0.5$ 对称。  
% 这种不对称性源于滑动侧链的引入打破了传统 AB 二嵌段共聚物的拓扑对称性，使体系在结构上类似于引入具有 A$_2$B 型特征的不对称链段，从而诱导出宏观相图的偏斜。  
% 在这种情况下，$\bar{\phi}_\mathrm{A}<0.5$ 与 $\bar{\phi}_\mathrm{A}>0.5$ 分别对应滑动嵌段主要位于胶束核心（core）与壳层（shell）时的相行为。  
% 因此，相图在低组分一侧（$\bar{\phi}_\mathrm{A}<0.5$）整体收缩，而在高组分一侧（$\bar{\phi}_\mathrm{A}>0.5$）则相应扩展，这一特征在高 $N\chi$ 区域尤为明显，例如高 $\phi_\mathrm{A}$ 区域的 BCC 相区。
Another important feature of Fig.~\ref{fig:PhaseDiagram} is that, when $\phi_{\mathrm a}\neq 0$, the phase boundaries are no longer symmetric about $\bar{\phi}_{\mathrm A}=0.5$.
This asymmetry originates from the sliding side chain, which breaks the topological symmetry of AB diblock copolymers.
As a result, the sliding side chain introduces an A$_2$B-like architectural asymmetry, which leads to a skewed phase diagram at the macroscopic level.
Accordingly, relative to the linear AB diblock reference, the phase diagram is compressed for $\bar{\phi}_{\mathrm A}<0.5$ and expanded for $\bar{\phi}_{\mathrm A}>0.5$.
This asymmetry reflects two different packing environments of the sliding A segment: it is located predominantly in the micellar core for $\bar{\phi}_{\mathrm A}<0.5$, whereas it mainly contributes to the shell for $\bar{\phi}_{\mathrm A}>0.5$.
This feature is particularly evident in the high-$N\chi$ region, for example in the BCC phase region at high $\bar{\phi}_{\mathrm A}$.


\begin{figure*}[!ht]
    \centering
    \includegraphics[width=\linewidth]{./Figure/Gyroid.png}
    \caption{
        (a) Three-dimensional isosurfaces of the A- and B-component density distributions at $\phi_{\mathrm a}=0.50$, $\bar{\phi}_{\mathrm A}=0.30$, and $N\chi=30$.
        (b) Spatial distributions of the blocks in L--S TT copolymers for $\phi_{\mathrm a}=0.25$, $0.50$, and $0.75$.
        Rows show two-dimensional density slices along the main diagonal ($x=y$), corresponding to the light-gray section in panel (a), for A1, A2, A3, and B, respectively.
        The densities of all A blocks are normalized by their corresponding overall volume fractions, $\phi_{\mathrm a}$ or $(1-\phi_{\mathrm a})$.
        (c) Probability distributions of the sliding degree of freedom, $P_\Lambda(f)$, for different $\phi_{\mathrm a}$ values, as indicated by color.
    }
    \label{fig:Gyroid}
\end{figure*}


% 为进一步揭示滑动拓扑自由度如何在分子尺度上响应上述相图变化，我们在 Gyroid 结构中系统分析各链段及拓扑参数 $f$ 的空间分布，从而刻画滑动侧链在三维有序结构中的空间自适应调节机制。
% 如图~\ref{fig:Gyroid} 所示，在设定 Ia$\overline{3}$d 空间群对称性下，在确定的组分比 $\bar{\phi}_\mathrm{A} = 0.30$，相互作用参数 $N\chi = 30.0$ 的条件下，对比 $\phi_{\mathrm{a}} = 0.25$, $0.50$, $0.75$ 的三种可滑动侧链体积分数的 L-S TT，均形成由可滑动链构成胶束核心的 Gyroid 相结构。
% 图~\ref{fig:Gyroid}a 展示了以 $\phi_{\mathrm{a}}=0.50$ 为例的典型形貌的三维视图。
% 图~\ref{fig:Gyroid}b 展示了在 $\phi_\mathrm{a}=0.25, 0.50, 0.75$ 三种不同侧链体积分数下，L–S TT 共聚物在 $(110)$ 晶面上的二维截面（见图~\ref{fig:Gyroid}a 中的灰色观测面）上的各嵌段 $m$ 的密度分布 $\phi_m(\bm r)$。
% 这些链段密度分布可由公式~(\ref{eq:BlockDensity}) 计算得到。
% 为便于不同 $\phi_\mathrm{a}$ 条件之间的对比，各 A 组分链段的密度已进行标准化处理，分别表示为 $\phi_\mathrm{A1}/\phi_\mathrm{a}$、$\phi_\mathrm{A2}/(1-\phi_\mathrm{a})$、$\phi_\mathrm{A3}/(1-\phi_\mathrm{a})$。  
% 图中每列对应不同的 $\phi_\mathrm{a}$，每行依次表示滑动侧链 A1、主链中的 tail 链段 A2、trunk 链段 A3 以及主链上的 B 组分的空间分布形态。
% 一般地，滑动侧链 $\phi_\mathrm{A1}$ 主要集中在 A 域的核心位置，形成致密的 A 核心区；trunk 链段靠近分子的接枝点，$\phi_\mathrm{A3}$ 因而位于 A/B 界面附近，包覆在 A 核外层。
% 此外，在不同滑动比例下 $\phi_\mathrm{B}$ 的空间分布形态基本一致，始终填充于 A 相胶束之间的间隙，表明不同 $\phi_\mathrm{a}$ 条件下胶束尺寸基本相同。
% 因此，不同 $\phi_\mathrm{a}$ 的差异主要体现在胶束内部各链段的分布上：
% 如 $\phi_\mathrm{a} = 0.25$ 时，$\phi_{\mathrm{A1}}$ 与 $\phi_{\mathrm{A2}}$都集中在 A 域的核心区域，对应于 RPA 预测的 “Y” 型 $\mathrm{A}_2\mathrm{AB}$ 构象，其中 tail 嵌段长度与滑动侧链长度相当\cite{Zhao2024Phase}。
% 而随 $\phi_\mathrm{a}$ 的增大，$\phi_\mathrm{A2}$ 的空间分布密度峰值逐渐由 A 核心区域向 A/B 界面外移。
% 因此在$\phi_\mathrm{a} = 0.75$的密度热力图中，$\phi_\mathrm{A1}$、$\phi_\mathrm{A2}$ 和 $\phi_\mathrm{A3}$ 呈现出由内至外依次嵌套的空间分布。
% 这种变化趋势表明，随着 $\phi_\mathrm{a}$ 的增大，L–S TT 分子趋向于形成更接近 AB 线型 diblock 的构型，以降低整体自由能。
% 对应的，$\phi_\mathrm{A1}$的空间分布则越来越向 A 组分的中间区域聚集；
% 而与 B 组分直接连接的 $\phi_\mathrm{A3}$ 始终分布于 A 组分区域的外壳层，且其空间位置基本不随 $\phi_\mathrm{a}$ 变化而改变，显示出其分布的相对稳定性。
To resolve the molecular origin of these phase diagram changes, we analyze the spatial distributions of the four blocks $m$ and the topology distribution $P_\Lambda(f)$ in the Gyroid phase.
This allows us to characterize the spatially adaptive regulation of the sliding side chain.
As shown in Fig.~\ref{fig:Gyroid}, under the prescribed Ia$\overline{3}$d space-group symmetry, L--S TT copolymers with three representative sliding side-chain fractions, $\phi_{\mathrm a}=0.25$, $0.50$, and $0.75$, all form a Gyroid phase with the micellar core composed of the sliding A component at $\bar{\phi}_{\mathrm A}=0.30$ and $N\chi=30.0$.
Fig.~\ref{fig:Gyroid}a shows a representative three-dimensional morphology for $\phi_{\mathrm a}=0.50$.
Fig.~\ref{fig:Gyroid}b shows two-dimensional density slices on the $(110)$ plane, corresponding to the gray observation plane in Fig.~\ref{fig:Gyroid}a, for the density distributions $\phi_m(\bm r)$ of each block $m$ at $\phi_{\mathrm a}=0.25$, $0.50$, and $0.75$.
These block density distributions are calculated from Eq.~(\ref{eq:BlockDensity}).
For comparison of the spatial block-density distributions at different $\phi_{\mathrm a}$, we rescale $\phi_{\mathrm{A1}}$ by $\phi_{\mathrm a}$ and both $\phi_{\mathrm{A2}}$ and $\phi_{\mathrm{A3}}$ by $1-\phi_{\mathrm a}$, corresponding to the prescribed fractions of A segments in the side-chain block and the main-chain A block, respectively.
This reduces composition driven differences in the overall density levels while preserving the spatial morphology and the relative distribution of A2 and A3.

In Fig.~\ref{fig:Gyroid}b, each column corresponds to a different $\phi_{\mathrm a}$, and the rows show, from top to bottom, the spatial distributions of the sliding side chain A1, the tail block A2, the trunk block A3, and the B block on the main chain.
In general, the sliding side chain $\phi_{\mathrm{A1}}$ is mainly concentrated in the core of the A domain, forming a dense A-rich core.
Because the trunk block is adjacent to the grafting point, $\phi_{\mathrm{A3}}$ is located near the A/B interface and forms an outer layer around the A core.
In addition, the spatial distribution of $\phi_{\mathrm B}$ remains nearly unchanged for different $\phi_{\mathrm a}$ and always fills the interstitial regions between A-rich micelles, indicating that the micellar size is almost unchanged as $\phi_{\mathrm a}$ varies.
Therefore, the main difference among different $\phi_{\mathrm a}$ lies in the internal partitioning of the A blocks within the micelles.
For $\phi_{\mathrm a}=0.25$, both $\phi_{\mathrm{A1}}$ and $\phi_{\mathrm{A2}}$ are concentrated in the core region of the A domain, corresponding to the ``Y''-shaped $\mathrm{A}_2\mathrm{AB}$ conformation predicted by RPA, where the tail block has a length comparable to that of the sliding side chain~\cite{Zhao2024Phase}.
As $\phi_{\mathrm a}$ increases, the density maximum of $\phi_{\mathrm{A2}}$ gradually shifts outward from the A-core region toward the A/B interface.
Consequently, for $\phi_{\mathrm a}=0.75$, the density maps show a nested spatial arrangement of $\phi_{\mathrm{A1}}$, $\phi_{\mathrm{A2}}$, and $\phi_{\mathrm{A3}}$ from the inner core to the outer interface.
This trend indicates that, with increasing $\phi_{\mathrm a}$, L--S TT molecules tend to adopt configurations closer to the linear AB diblock limit in order to lower the total free energy.
Correspondingly, the spatial distribution of $\phi_{\mathrm{A1}}$ becomes increasingly localized toward the central region of the A domain.
By contrast, $\phi_{\mathrm{A3}}$, which is directly connected to the B block, remains located in the outer shell of the A domain and shows little change in spatial position with varying $\phi_{\mathrm a}$.


% 对比不同 $\phi_\mathrm{a}$ 下 Gyroid 相的全空间滑动分布 $P_\Lambda(f)$ 的变化情况，可以进一步表征体系中滑动自由度带来的拓扑结构自适应特征。
% 固定  $N\chi = 30.0$，如图~\ref{fig:Gyroid}c，发现不同 $\phi_{\mathrm a}$ 条件下的 $P_{\Lambda}(f)$ 均呈现唯一的极大值，表明各拓扑 $f_j$ 并非等概率出现，而体系对某一特定拓扑具有明显偏好。
% 该峰值对应的拓扑定义为最优拓扑 $f^{*}$，并在图~\ref{fig:Gyroid}c 中以虚线标示。
% 随着可滑动链段比例 $\phi_\mathrm{a}$ 的增加，最优拓扑 $f_\Lambda^\ast$ 的取值逐渐降低。
% 这表明当可滑动侧链变长时，L-S TT 自适应形成类似 AB 二嵌段共聚物的线性拓扑结构，而非保持滑动链的 $\mathrm{A_2 AB}$ 星型拓扑结构，我们称其为线型拓扑择优。
% 这是 L-S TT 在自组装过程中，为了兼顾界面曲率匹配与填充挫折的缓解，
% 而自发调整滑动拓扑所形成的选择结果，与无规相近似理论（RPA）中的预测相符\cite{Zhao2024Phase}。
The variation of the full-space sliding distribution $P_\Lambda(f)$ for the Gyroid phase at different $\phi_{\mathrm a}$ further characterizes the adaptive topology selection induced by the sliding degree of freedom.
At fixed $N\chi=30.0$, as shown in Fig.~\ref{fig:Gyroid}c, $P_\Lambda(f)$ exhibits a single maximum for each $\phi_{\mathrm a}$, indicating that different topologies $f_j$ are not equally sampled and that the system has a preference for a specific topology.
The topology corresponding to this maximum is defined as the preferred topology $f_\Lambda^\ast$ and is marked by the dashed lines in Fig.~\ref{fig:Gyroid}c.
As the sliding chain fraction $\phi_{\mathrm a}$ increases, the value of $f_\Lambda^\ast$ gradually decreases.
This indicates that, when the sliding side chain becomes longer, L--S TT copolymers preferentially adopt a more linear AB-diblock-like topology rather than maintaining an $\mathrm{A_2AB}$ star-like topology, which we refer to as linear topology preference.
This preference results from the spontaneous adjustment of sliding topology during self-assembly, allowing the system to respond to variations in interfacial curvature and heterogeneous packing constraints, and is consistent with the prediction of the RPA~\cite{Zhao2024Phase}.

\PDFsubsection{Topology Response to Interfacial Curvature}
    In conventional block copolymers, spontaneous curvature is usually tuned by prescribed molecular parameters, such as the Kuhn length, component volume fraction, or grafting number\cite{Xu2024Architectural}.
In contrast, L--S TT copolymers can locally adjust the preferred topology $f_{\mathcal{R}}^\ast$ through sliding, because AB linear diblock and A$_2$B star topologies possess different curvature preferences.
As a weakly curved three-dimensional network phase whose interfacial curvature is closest to that of the lamellar phase, the $\mathrm{O}^{70}$ phase provides a useful model for examining how topology transformability affects interfacial adaptation.
We therefore use the $\mathrm{O}^{70}$ phases at low and high $\bar{\phi}_{\mathrm A}$ in the low-$N\chi$ region to examine how the sliding degree of freedom affects microphase separation and responds to interfacial curvature.


As shown in Fig.~\ref{fig:PhaseDiagram}, increasing $\phi_\mathrm{a}$ shifts the $\mathrm{O}^{70}$ phase region toward lower $\bar{\phi}_\mathrm{A}$ relative to the linear AB diblock reference in both the A-core-forming regime ($\bar{\phi}_\mathrm{A}<0.5$) and the A-shell-forming regime ($\bar{\phi}_\mathrm{A}>0.5$).
This trend originates from the reduction of the radius of gyration $R_g$ of the A block caused by the sliding side chain, which decreases the spontaneous curvature associated with the A component.
However, the shifts along the $N\chi$ direction differ between the A-core-forming and A-shell-forming regimes.
Beyond the slight increase in $N\chi_{\mathrm{ODT}}^{\mathrm{Slide}}$ relative to the linear AB diblock reference, the $\mathrm{O}^{70}$ phase region shifts to higher $N\chi$ in the A-core-forming regime but to lower $N\chi$ in the A-shell-forming regime.
This asymmetry reflects the opposite curvature requirements of the A component in the two regimes: positive spontaneous curvature when A forms the core and negative spontaneous curvature when A forms the shell regions.
Because the sliding side chain reduces the spontaneous curvature of the A block, the statistical weights of different $f$ respond differently to the curvature environments in the two regimes.
The A-core-forming regime is associated with a larger contribution from more negative-curvature topologies, whereas the A-shell-forming regime favors nearly linear configurations that are more compatible with its local geometry.
This curvature-conditioned topology response is consistent with the asymmetric displacement of the $\mathrm{O}^{70}$ phase region in the phase diagram.


\begin{figure}[!htb]
    \centering
    \includegraphics[]{./Figure/O70-Pf.png}
    \caption{
        (a,b) Three-dimensional views of the $\mathrm{O}^{70}$ phase formed by L--S TT copolymers at $\bar{\phi}_{\mathrm A}=0.42$ and $0.55$, respectively.
        Red--yellow regions denote the A-component density.
        (c,d) Isodensity-layer topology distributions $P_\Omega(f)$ at $\phi_{\mathrm A}^{\mathrm{iso}}=0.75$ as functions of the sliding side-chain fraction $\phi_{\mathrm a}$ for the two compositions in panels (a) and (b), respectively.
        Heat maps show $P_\Omega(f)$, contours indicate numerical isovalues, and the red and blue lines denote the locations of the maximum and minimum of the distribution, $f_\Omega^\ast$ and $f_\Omega^-$, with symbols showing the corresponding numerical results.
        (e,f) Isodensity-layer topology distributions $P_\Omega(f)$ as functions of the sampled density level $\phi_{\mathrm A}^{\mathrm{iso}}$ at fixed $\phi_{\mathrm a}=0.50$, for $\bar{\phi}_{\mathrm A}=0.42$ and $0.55$, respectively, with the same notation as in panels (c,d).
        (g,h) Schematic illustrations of the topology redistribution of sliding chains across the interfacial region for $\bar{\phi}_{\mathrm A}=0.42$ and $0.55$, respectively.
    }
    \label{fig:O70Topology}
\end{figure}


To examine the topology response in the O$^{70}$ phase, we compare the Core ($\bar{\phi}_{\mathrm A}=0.42$) and Shell ($\bar{\phi}_{\mathrm A}=0.55$) states at $N\chi=13.0$, both with Fddd symmetry (Fig.~\ref{fig:O70Topology}a,b).
Fig.~\ref{fig:O70Topology}c,d shows their isodensity layer topology distributions $P_\Omega(f)$ at $\phi_{\mathrm A}^{\mathrm{iso}}=0.75$ as $\phi_{\mathrm a}$ varies.
At this A-rich density level, $P_\Omega(f)$ remains single peaked over the examined $\phi_{\mathrm a}$ range, with the red line tracing the preferred topology $f_\Omega^\ast$.
At small $\phi_{\mathrm a}$, $f_\Omega^\ast$ remains near $0.3$, corresponding to a typical ``Y''-shaped $\mathrm{A}_2\mathrm{AB}$ topology, and then decreases toward the AB-like limit as $\phi_{\mathrm a}$ increases.
This ordered-phase topology preference is consistent with the RPA prediction near the ODT and motivates the following analysis of its spatial relation to interfacial geometry.


To resolve the spatial dependence of topology selection, we examine $P_\Omega(f)$ across isodensity layers characterized by $\phi_{\mathrm A}^{\mathrm{iso}}$ at fixed $\phi_{\mathrm a}=0.50$.
As shown in Fig.~\ref{fig:O70Topology}e,f, the sampled layers span the A-rich, interfacial, and B-rich regions as $\phi_{\mathrm A}^{\mathrm{iso}}$ decreases from $0.90$ to $0.10$.
In the A-rich region, $P_\Omega(f)$ is single peaked, with its maximum at $f_\Omega^\ast$, whereas toward the B-rich region it develops a U-shaped profile with enhanced weights near the AB-like ($f=0$) and A$_2$B-like ($f=1$) limits and a local minimum $f_\Omega^-$ between them.
For $\phi_{\mathrm{A}}^{\mathrm{iso}} \lesssim 0.40$, the endpoint contributions become increasingly pronounced while intermediate sliding configurations remain populated.
This layer dependence reveals a spatial redistribution of topology weights across the interface, as illustrated schematically in Fig.~\ref{fig:O70Topology}g,h.
The A-segment population in the A-rich region is preferentially contributed by A$_2$AB-like configurations, whereas toward the B-rich region the relative contributions shift toward the AB-like and A$_2$B-like limits.
This redistribution is consistent with their different natural extension lengths: the more compact A$_2$B-like configurations preferentially contribute near the interface, whereas the more extended AB-like configurations can reach farther toward the A-rich region while their B blocks remain closer to the B-rich domain interior.
The enhanced A$_2$AB-like contribution in the A-rich region reflects the participation of its two A branches in filling that domain.
The sliding degree of freedom therefore redistributes statistical weight within a single continuous topology ensemble across different spatial regions, rather than producing distinct topological populations; this interpretation is supported by the weak spatial variation of $\bar f_{\mathcal R}$ and the small $D_{\mathcal R}$ discussed in Sec.~\ntxref{SI-subsec:TopologyStatisticsSpatial} of the Supporting Information.



% The response of $P_\Omega(f)$ to $\phi_{\mathrm A}^{\mathrm{iso}}$ further demonstrates the interfacial-curvature adaptability unique to L--S TT copolymers.
% This adaptability manifests differently in the two $\mathrm{O}^{70}$ phases because the A component plays distinct structural roles in the A-core-forming and A-shell-forming regimes.
% For $\bar{\phi}_{\mathrm A}=0.42$, where A forms the network core, $f_\Omega^\ast$ jumps from the A$_2$B endpoint to the AB linear diblock endpoint near $\phi_{\mathrm A}^{\mathrm{iso}}\approx0.23$, stays close to zero over an intermediate range, and then increases toward an A$_2$AB-like topology at larger $\phi_{\mathrm A}^{\mathrm{iso}}$.
% This behavior can be attributed to the larger spatial volume occupied by the two A branches in the A$_2$B topology, which provides a lower, or even negative, A-side spontaneous curvature.
% Such a topology can compensate for the increased positive curvature associated with $\bar{\phi}_{\mathrm A}<0.5$ at the interface, thereby satisfying the curvature requirement of the $\mathrm{O}^{70}$ phase.
% In the A-core-forming geometry, the outer region of the A core requires additional A segments to maintain density continuity.
% As the interfacial stretching penalty increases, the preferred configuration shifts from an AB linear diblock topology toward a Y-shaped A$_2$AB topology, filling the interfacial region more efficiently and lowering the free energy.
% By contrast, for $\bar{\phi}_{\mathrm A}=0.55$, the low-$\phi_{\mathrm A}^{\mathrm{iso}}$ region contains a larger contribution from AB linear diblock configurations, with $f^*\approx0$.
% This is because the A-shell-forming geometry already favors lower, or negative, spontaneous curvature, so the system can maintain a low interfacial energy without converting strongly toward A$_2$B-like topologies.
% Consequently, the onset of $f^*>0$ is delayed to around $\phi_{\mathrm A}^{\mathrm{iso}}\approx0.50$, and the variation of $f_\Omega^\ast$ becomes smoother than that for $\bar{\phi}_{\mathrm A}=0.42$, indicating a more stable curvature-driven conformational adjustment.
The response of $P_\Omega(f)$ to $\phi_{\mathrm A}^{\mathrm{iso}}$ further reveals the interfacial-curvature response of L--S TT copolymers.
This response differs between the two O$^{70}$ states because the A component plays distinct structural roles in the A-core-forming and A-shell-forming geometries.
For $\bar{\phi}_{\mathrm A}=0.42$, where A forms the network core, $f_\Omega^\ast$ shifts from the A$_2$B-like endpoint to the AB-like limit near $\phi_{\mathrm A}^{\mathrm{iso}}\approx0.23$, remains close to zero over an intermediate range, and then increases toward an A$_2$AB-like topology at higher $\phi_{\mathrm A}^{\mathrm{iso}}$.
This behavior is consistent with the larger spatial extent of the two A branches in the A$_2$B-like topology, which favors a lower, or more negative, A-side spontaneous curvature.
This lower-curvature tendency can compensate for the more positive curvature required by the A-core geometry.
As the sampled layer moves inward toward the A-rich core, the preferred topology shifts from the AB-like limit toward A$_2$AB-like configurations, whose two A branches contribute more strongly to filling the A-rich region.
By contrast, for $\bar{\phi}_{\mathrm A}=0.55$, where A forms the shell, the low-$\phi_{\mathrm A}^{\mathrm{iso}}$ region retains a larger AB-like contribution with $f_\Omega^\ast\approx0$.
Because the A-shell geometry favors lower, or more negative, A-side curvature, the topology redistribution required to accommodate the local geometry differs from that in the A-core state.
Consequently, the onset of finite $f_\Omega^\ast$ occurs at a higher $\phi_{\mathrm A}^{\mathrm{iso}}$, around $0.50$, and its variation with density level is smoother than in the Core state.
These contrasting responses therefore reflect the different curvature requirements of the A-core and A-shell geometries and the corresponding redistribution among accessible topologies with different curvature preferences.
Because $P_\Omega(f)$ characterizes this layer dependent topology redistribution but does not directly quantify the local interfacial geometry, we further examine the relation between topology selection and local curvature.


To characterize the interfacial geometry directly, we calculate the local mean curvature from the A-component density field.
For an A-component isodensity surface, the unit normal is defined as $\bm n(\bm r)=-\nabla\phi_{\mathrm A}(\bm r)/|\nabla\phi_{\mathrm A}(\bm r)|$, pointing from the A-rich region toward the B-rich region.
The local mean curvature is then evaluated as~\cite{Cheng2024Understand}
\begin{equation}
    H(\bm r) = \frac{1}{2}\nabla\cdot\bm n(\bm r).
\end{equation}
With this convention, positive $H$ corresponds locally to an A-core-like curvature, whereas negative $H$ corresponds to an A-shell-like curvature.
In the numerical evaluation, the spatial gradients and divergence are computed using periodic spectral derivatives, with a weak Gaussian smoothing applied to $\phi_{\mathrm A}(\bm r)$ to suppress grid-scale fluctuations in the curvature.


\begin{figure}[!htb]
    \centering
    \includegraphics[]{./Figure/O70-Curvature.png}
    \caption{
        (a,b) Mean curvature fields $H (R^{-1}_{\mathrm g})$ on the central $z$ section for the Core ($\bar{\phi}_{\mathrm A}=0.42$) and Shell ($\bar{\phi}_{\mathrm A}=0.55$) states; cyan and magenta contours mark $\phi_{\mathrm A}^{\mathrm{iso}}=0.50$ and $0.80$.
        (c,d) $f_{\mathcal R}^{\ast}(\bm r)$ versus $H (R^{-1}_{\mathrm g})$ on the corresponding three-dimensional isodensity surfaces at $\phi_{\mathrm A}^{\mathrm{iso}}=0.50$ and $0.80$, respectively.
        Dark red and orange symbols denote the Core ($\bar{\phi}_{\mathrm A}=0.42$) and Shell ($\bar{\phi}_{\mathrm A}=0.55$) states, respectively.
    }
    \label{fig:O70Curvature}
\end{figure}


Fig.~\ref{fig:O70Curvature}a,b shows the local mean-curvature fields together with the $\phi_{\mathrm A}^{\mathrm{iso}}=0.50$ and $0.80$ isodensity contours.
The former samples the A/B interfacial region, whereas the latter probes a more A-rich region.
For Fig.~\ref{fig:O70Curvature}c,d, spatial points $\bm r$ are sampled over the entire three-dimensional isodensity surface to obtain $H(\bm r)$ and $P_{\mathcal R(\bm r)}(f)$, from which the corresponding $f_{\mathcal R}^{\ast}(\bm r)$ values are determined.
A clear distinction between the Core and Shell states is observed at the two density levels.
On the interfacial surface, $\phi_{\mathrm A}^{\mathrm{iso}}=0.50$, the Core state is biased toward more positive $H$ and exhibits larger $f_{\mathcal R}^{\ast}$, whereas the Shell state is biased toward more negative $H$ and smaller $f_{\mathcal R}^{\ast}$.
On the more A-rich surface, $\phi_{\mathrm A}^{\mathrm{iso}}=0.80$, the relative ordering of $f_{\mathcal R}^{\ast}$ between the two states is reversed, highlighting the distinct topology responses associated with the A-core and A-shell geometries.
Despite these differences, $f_{\mathcal R}^{\ast}$ shows a weak overall tendency to increase with increasing $H$ within the sampled surfaces, although appreciable scatter remains and the trend becomes less distinct near $H\approx0$.
These results show that the topology--curvature response is strongly conditioned by whether A forms the network core or the surrounding shell, while retaining an overall sensitivity of the preferred topology to the local mean curvature.



Taken together, the asymmetric displacement of the O$^{70}$ phase region reflects the coupling between sliding topology redistribution and the local curvature environment.
When A forms the core, topology selection responds to a more positive curvature environment, whereas when A forms the shell, it is redistributed to accommodate a lower, or more negative, curvature requirement.
The analysis of $H$ and $f_{\mathcal R}^{\ast}$ further shows that the two geometries exhibit distinct responses of topology to curvature, while the preferred topology retains a weaker overall dependence on the local mean curvature.
At larger $\phi_{\mathrm a}=0.75$, topology locking reduces the available curvature adjustment, and the high-$\bar{\phi}_{\mathrm A}$ O$^{70}$ region correspondingly shifts back toward higher $N\chi$ (Fig.~\ref{fig:PhaseDiagram}c).

\PDFsubsection{Topology Response to Heterogeneous Packing Environments}
    Packing frustration refers to the inability of block copolymers to simultaneously produce uniform interfacial curvature for minimizing interfacial energy and domains of uniform thickness for minimizing chain-stretching energy in nonlamellar structures~\cite{Matsen1997Block,Xu2024Architectural}.
A common mechanism for mitigating packing frustration is local segregation, where differences in block type, composition, architecture, or Kuhn length allow different chain blocks to occupy different regions of a micelle, thereby reducing local stretching inhomogeneity and curvature-induced strain~\cite{Xu2024Architectural}.


L--S TT copolymers provide a topology-based route for responding to heterogeneous local filling environments associated with packing frustration by continuously adjusting among AB, A$_2$AB, and A$_2$B configurations through the sliding degree of freedom.
This effect is particularly relevant in the spherical BCC and FCC phases, where the Wigner--Seitz cells contain heterogeneous local filling environments around each B-rich domain.
At high $N\chi$, stronger segregation restricts chain conformations and amplifies the mismatch between local packing geometry and the characteristic chain extension favored by different regions of the cell.
The sliding topology can respond to these heterogeneous constraints by redistributing its statistical weight according to the local environment.
We therefore examine this adaptive topology redistribution in the strongly segregated BCC and FCC phases.


\begin{figure}[!t]
    \centering
    \includegraphics[]{./Figure/Sphere-fstar.png}
    \caption{
        (a,b) Density distributions on the [110] section, $x=y$, for representative BCC (a, $\bar{\phi}_{\mathrm A}=0.85$) and FCC (b, $\bar{\phi}_{\mathrm A}=0.885$) phases at $N\chi=45.0$ and $\phi_{\mathrm a}=0.50$.
        The upper and lower heat maps show $\phi_{\mathrm{A1}}(\bm r)$ and $\phi_{\mathrm{A2}}(\bm r)$, respectively; white dashed lines mark isodensity contours.
        (c,d) Local A density $\phi_{\mathrm A}(\bm r)$ (c) and preferred topology $f_{\mathcal R}^{\ast}(\bm r)$ (d) resolved by the distance $l_{\mathrm B}$ to the nearest B-domain center.
        Colored violins denote the BCC and FCC distributions within each $l_{\mathrm B}$ interval, with gray points showing sampled local values; only A-rich positions with $\phi_{\mathrm A}>0.5$ are included in (d).
        (e) Spatial distributions of $f_{\mathcal R}^{\ast}(\bm r)$ on the BCC (left) and FCC (right) Wigner--Seitz cell sections, shown with a common color scale.
        Blue spheres denote the centers of the B-component micelles, and gray frames denote the Wigner--Seitz cells.
    }
    \label{fig:SphereFStar}
\end{figure}


At $N\chi=45.0$, the systems with $\bar{\phi}_{\mathrm A}=0.85$ and $0.885$ form the BCC and FCC phases, respectively.
For $\phi_{\mathrm a}=0.50$, the A1 sliding block and the A2 tail exhibit distinct spatial distributions in both BCC and FCC phases (Fig.~\ref{fig:SphereFStar}a,b), despite their otherwise interchangeable roles in the absence of sliding.
To relate this topology response to the local geometry, we characterize each spatial position by its distance $l_{\mathrm B}$ from the nearest B-domain center.
Specifically, we define
\begin{equation}
    l_{\mathrm B}(\bm r)
    =
    \min_j
    \left|
        \bm r-\bm r_{\mathrm B,j}
    \right|,
\end{equation}
where $\bm r_{\mathrm B,j}$ denotes the center of the $j$th B-rich domain, and the distance is evaluated under periodic boundary conditions.
As $l_{\mathrm B}/R_{\mathrm g}$ increases beyond approximately $1.2$, the sampled positions enter the A-rich interior and $\phi_{\mathrm A}(\bm r)$ approaches its high-density regime, whereas $f_{\mathcal R}^{\ast}(\bm r)$ continues to span a broad range of values (Fig.~\ref{fig:SphereFStar}c,d).
Positions with similar $l_{\mathrm B}$ are not necessarily locally equivalent, because their orientations and neighboring packing environments within the Wigner--Seitz cell can differ.
Accordingly, $f_{\mathcal R}^{\ast}$ cannot be characterized by $l_{\mathrm B}$ alone, but depends on the specific local geometry.
The Wigner--Seitz cell maps in Fig.~\ref{fig:SphereFStar}e make this dependence on local geometry explicit.
Even at comparable $l_{\mathrm B}$ within the A-rich region, different positions on the Wigner--Seitz cell section exhibit distinct $f_{\mathcal R}^{\ast}$ values, indicating that the preferred topology depends on the local packing environment rather than on the nearest-domain distance alone.
The preferred topology is therefore determined by the specific local packing environment rather than by a single radial distance, consistent with the direction-dependent responses discussed further in Supporting Information Sec.~\ntxref{SI-sec:TopologyHomogenization}.


\begin{figure}[!t]
    \centering
    \includegraphics[]{./Figure/Sphere-POmega.png}
    \caption{
        (a,b) Isodensity layer topology distributions $P_{\Omega}(f)$ at $\phi_{\mathrm A}^{\mathrm{iso}}=0.75$ for $\phi_{\mathrm a}=0.25$, $0.50$, and $0.75$ in the BCC (a) and FCC (b) phases at $N\chi=45$.
        Stars mark the preferred topology $f_{\Omega}^{\ast}$, with dashed lines indicating the corresponding values of $f$.
        (c) Distributions of the local KL divergence $D_{\mathcal R}(\bm r)$ in the A-rich regions, $\phi_{\mathrm A}(\bm r)>0.5$, for representative Gyroid, BCC, and FCC states at $N\chi=45$ and $\phi_{\mathrm a}=0.50$.
        Red and orange violins denote the Core and Shell states, corresponding to the lower and higher $\bar{\phi}_{\mathrm A}$ branches of each morphology, respectively.
        For G, BCC, and FCC, the corresponding Core/Shell pairs of $\bar{\phi}_{\mathrm A}$ are $(0.30,\,0.54)$, $(0.16,\,0.85)$, and $(0.14,\,0.885)$, respectively.
        Gray points show sampled local values, and circles denote the corresponding sliding distribution entropy loss $\Delta S_{\mathrm{Slide}}/k_{\mathrm B}$.
    }
    \label{fig:SpherePOmega}
\end{figure}

The pronounced spatial variation of $f_{\mathcal R}^{\ast}$ does not imply a stronger localization of the full topology distribution.
At individual spatial positions, the local distributions $P_{\mathcal R}(f)$ already remain broad and only weakly localized around their preferred topologies.
At the same time, $f_{\mathcal R}^{\ast}$ varies among different local environments, so that spatial averaging over an isodensity layer combines local distributions centered at different preferred values of $f$.
We refer to the combined effect of these two contributions as topology homogenization: the broad local accessibility of sliding configurations, together with the spatial redistribution of their preferred values, produces an overall topology distribution that is close to uniform sliding.

Compared with the single-peaked distribution in Fig.~\ref{fig:Gyroid}c, the isodensity-layer distributions $P_{\Omega}(f)$ in the spherical phases remain broad over the accessible range of $f$ for different $\phi_{\mathrm a}$ (Fig.~\ref{fig:SpherePOmega}a,b).
In the FCC phase, the separation between different local topology preferences is sufficiently pronounced that $P_{\Omega}(f)$ even develops a bimodal profile for some $\phi_{\mathrm a}$.
The representative local distributions in Supporting Information Sec.~\ntxref{SI-sec:TopologyHomogenization} further show that different local environments can shift $f_{\mathcal R}^{\ast}$ while the $P_{\mathcal R}(f)$ remain broad.
This provides a statistical basis for topology homogenization in the spherical phases.


The degree of topology homogenization can be quantified by the local KL divergence $D_{\mathcal R}(\bm r)$ and the corresponding spatial average $\Delta S_{\mathrm{Slide}}$ weighted by the local A-segment population.
They measure the local and overall departures from uniform sliding, respectively, as defined in Eqs.~\eqref{eq:LocalDivergence} and \eqref{eq:SlideEntropy}.
For the representative states in Fig.~\ref{fig:SpherePOmega}c, $\Delta S_{\mathrm{Slide}}$ follows the same overall ordering, G $>$ BCC $>$ FCC, whether the A component forms the core or the shell.
The persistence of this trend across the two structural roles of A indicates that the degree of topology homogenization is primarily associated with the different local filling constraints of these structures.

Within each phase, $\Delta S_{\mathrm{Slide}}$ is consistently slightly smaller when A forms the shell than when it forms the core, indicating a modest curvature effect on topology homogenization.
This effect is reflected more strongly in the width of the $D_{\mathcal R}$ distribution than in the value of $\Delta S_{\mathrm{Slide}}$ itself.
This interpretation is further supported by the spatial distributions of $D_{\mathcal R}$ in Supporting Information Fig.~\ntxref{SI-fig:PhaseTopology} and by the discussion in Supporting Information Sec.~\ntxref{SI-sec:TopologyHomogenization}.
Together with the directional topology distributions discussed in Supporting Information Sec.~\ntxref{SI-sec:TopologyHomogenization}, this homogenization suggests a microscopic picture in which the main chain primarily accommodates the connection between neighboring B-rich domains, while the sliding A1 branch adjusts its spatial contribution to fill the intervening regions.
The relative contributions of different topologies to the local A-segment density can therefore vary across different local filling environments without strongly restricting the accessible topology ensemble.


Thus, strong spatial variation of the preferred topology can coexist with a broadly accessible local sliding ensemble.
This topology-based response enables the system to adapt to heterogeneous local packing environments by redistributing statistical weight toward suitable topologies while retaining broad configurational accessibility, and may thereby help reduce the stretching inhomogeneity associated with packing frustration.


\PDFsubsection{Phase Stability under Enthalpy--Entropy Competition}
    \begin{figure*}[!htb]
    \centering
    \includegraphics[width=\linewidth]{./Figure/PhaseWidth.png}
    \caption{
        (a) Phase width variation for different sliding side-chain fractions, $\phi_{\mathrm a}=0.25$, $0.50$, and $0.75$.
        The upper line plot compares the normalized widths of the L, G, C, BCC, and FCC regions at $N\chi=40.0$, defined as $\Delta\bar{\phi}_{\mathrm A}(\phi_{\mathrm a})/\Delta\bar{\phi}_{\mathrm A}(0.25)$.
        The lower stacked bars show the phase widths along the $\bar{\phi}_{\mathrm A}$ direction for different $\phi_{\mathrm a}$ at $N\chi=25.0$, $35.0$, and $45.0$.
        (b--d) BCC phase region.
        (b) Width of the BCC region, $\Delta\bar{\phi}_{\mathrm A}$, as a function of $N\chi$ for $\phi_{\mathrm a}=0.25$, $0.50$, and $0.75$.
        C and S denote the A-core and A-shell composition ranges, respectively; circles and triangles denote the corresponding ranges, and dashed lines are spline fits.
        (c) Isodensity layer preferred topology $f_{\Omega}^{\ast}$ as a function of $\phi_{\mathrm A}^{\mathrm{iso}}$ at different $N\chi$ for fixed $\phi_{\mathrm a}=0.50$.
        (d) $f_{\Omega}^{\ast}$ as a function of $\phi_{\mathrm a}$ at fixed $\phi_{\mathrm A}^{\mathrm{iso}}=0.75$ for different $N\chi$.
        (e--g) Same as (b--d), but for the FCC phase.
    }
    \label{fig:PhaseWidth}
\end{figure*}


% 综上所述，滑动自由度 $f$ 提供了额外的构象调节机制，使 L-S TT 能够在不同空间位置自发调整拓扑，以调整自发曲率和适应有序相中的局部密度要求，缓解 packing frustration。
% 从自由能角度看，$f$ 的作用主要体现在两个方面。  
% 其一，$f$ 的引入扩展了体系可访问的拓扑态空间，从而带来由滑动自由度产生的额外熵贡献。  
% 其二，$f$ 的空间重分布会改变链在界面附近的伸展方式，使链具有了对局域密度匹配以及对曲率的适应能力。
% 这些由 $f$ 引入所带来的影响，最终都体现在自由能中熵项与焓项的竞争之中。
% 随着 $N\chi$ 增大，界面相关的焓贡献逐渐增强，因此滑动熵与局域焓优化这两种效应的相对权重也会发生变化，进而使 L-S TT 在低 $N\chi$ 下表现出拓扑择优，对应低滑动熵，而在高 $N\chi$ 下转而表现出拓扑均化，对应高滑动熵。
% 具体地，在低 $N\chi$ 区域，体系对构象调节的需求相对较弱，此时仅出现层状相及双联通相等界面形状较为简单、曲率变化较缓的结构，局域环境对拓扑的要求也相对一致，因而体系倾向于围绕某一择优的拓扑 $f_\Omega^\ast$ 形成单峰分布。
% 在这一条件下，SCFT 得到的最优拓扑与 RPA 单模近似预测的 $f_\Omega^\ast$ 近似一致，拓扑涨落所对应的额外构象熵主要体现在单峰分布宽度的变化上。
% 当 $N \chi$ 增大且体系形成的有序结构具有复杂界面（如 BCC 与 FCC 相）时，调节 $f$ 能显著降低界面能与链段伸展挫折所带来的能量代价。
% 此时，体系不再偏向单一的择优拓扑，而是呈现出更宽的 $f$ 分布范围；不同拓扑构型在魏格纳–赛茨元胞的不同空间位置上发生功能分化，并协同分担局域填充任务。
% 另一方面，即便在高 $N \chi$ 条件下，若目标相结构界面形貌较为规整（例如图~\ref{fig:Gyroid} 中的 Gyroid 相），调节 $f$ 带来的界面能优化有限，系统仍可能维持受构象熵主导的$f$分布。
In summary, the sliding degree of freedom $f$ provides an additional mechanism for conformational regulation, allowing L--S TT copolymers to redistribute the contributions of different topologies to the local A-segment density in response to local interfacial curvature and filling environments.
This local topological adaptation helps accommodate structural requirements in ordered phases and may thereby relieve packing frustration.


Allowing $f$ to vary expands the accessible molecular configurations and enables the topology weights to redistribute self-consistently in space.
The resulting topology redistribution modifies the chain conformational response near interfaces, coupling the conformational and sliding distribution entropy responses to local density matching and curvature response through the self-consistent fields.



As $N\chi$ increases, stronger A--B repulsion enhances the interfacial and packing constraints, changing the balance between conformational regulation and local geometric adaptation.

The observed evolution from pronounced topology selection at low $N\chi$ toward broader topology distributions at high $N\chi$ reflects this changing balance.

Specifically, in the low-$N\chi$ regime, the need for conformational adjustment is relatively weak.
Consistently, the stable ordered phases in this regime are mainly those with relatively simple interfacial geometry and low interfacial curvature, such as L and G phases.
Because the local environments impose similar topological requirements, the system tends to form a single-peaked distribution centered around a preferred topology $f_\Lambda^\ast$.

Under these conditions, the preferred topology obtained from SCFT is approximately consistent with $f_{\mathrm{RPA\text{-}MUT}}^\ast$ predicted by RPA--MUT.
The width of the single-peaked distribution characterizes the range of contributing topologies, while the associated conformational response is contained in the chain statistics.

When $N\chi$ increases and the ordered structures develop more complex interfaces, such as in the BCC and FCC phases, adjusting $f$ provides an additional route for accommodating interfacial geometry and heterogeneous local filling requirements.
The system then no longer favors a single preferred topology, but instead develops a broader distribution of $f$.
Different topological configurations become functionally differentiated at different positions within the Wigner--Seitz cell and cooperatively accommodate local filling requirements.

Even at high $N\chi$, an ordered structure with relatively regular interfacial geometry, such as the Gyroid phase in Fig.~\ref{fig:Gyroid}, may retain a pronounced topology preference.
This behavior suggests that the local geometric requirements favor similar topological configurations across different spatial regions.



% 通过比较不同 $N\chi$ 条件下相区宽度的变化，可以揭示体系中熵效应与焓效应之间的竞争关系。
% 图~\ref{fig:PhaseWidth}a 展示了在不同 $\phi_\mathrm{a}$ 下，L-S TT 体系随 $N\chi$ 增加所形成的有序相区宽度的变化趋势。
% 柱状图给出了不同相在总有序区宽度中的占比变化；上方折线图则标准化比较了各 $\phi_\mathrm{a}$ 条件下有序区宽度 $\Delta \bar{\phi}_{\mathrm{A}}(\phi_\mathrm{a},X)$ 的相对变化。
% 其中，$X \in \{\mathrm{L},\mathrm{G},\mathrm{C},\mathrm{BCC},\mathrm{FCC}\}$ 表示不同的有序相结构。
% 整体来看，随着 $N\chi$ 的增加，总有序相区的宽度呈增长趋势，符合嵌段共聚物在组分不相容性增强下，自发形成有序结构以降低焓代价的普遍特征。
% 然而，该增长趋势在不同 $\phi_\mathrm{a}$ 条件下的表现并不一致，这主要源于两方面机制的耦合：构象熵对 $\phi_\mathrm{a}$ 依赖的非线性特征与不同相所对应界面结构复杂度的差异。
% 首先，图~\ref{fig:ODT}c 已表明，L-S TT 对构象熵的调节能力并不随 $\phi_\mathrm{a}$ 单调增强，而是呈现出先增后降的非单调变化\fc{这里把长侧链和短侧链的论述删掉了，前文（第四节开始第一段，以及对滑动熵部分的讨论）已经讨论过这一结果，这里应该不用再重复叙述一遍。}。
% 相应地，在中等滑动比例 $\phi_\mathrm{a}\approx 0.50$ 附近，体系对滑动自由度的利用最为充分，因而其构象熵调节能力也最强。
% 其次，各有序相的界面结构复杂程度不同，对滑动调节能力的依赖也各异，从而导致不同相在 $\phi_\mathrm{a}$ 增加过程中的响应趋势有所不同。
% L 相具有规则平直界面，但要求 A、B 组分的体积比接近，而 L-S TT 一般不满足该条件，构象调节对其稳定性的增益有限，故其占比表现出先减后增的非单调行为，如图~\ref{fig:PhaseWidth}a 所示。
% G 相包含互穿通道，对链段的柔性调节更敏感，因此其相区占比随 $\phi_\mathrm{a}$ 增大而显著下降。
% 相比之下，C 相以及 BCC 与 FCC 等球状相结构复杂、界面高度弯曲且填充紧密，对滑动自由度调控密度场的能力依赖更强，因而其相区宽度在 $\phi_\mathrm{a} \approx 0.50$ 附近达到最大值。
% 对 $\phi_\mathrm{a}$ 依赖的差异性在相区宽度增幅比中同样存在，比较图~\ref{fig:PhaseWidth}a 中 $\phi_{\mathrm{a}}=0.25$ 和 $\phi_\mathrm{a}=0.50$ 的结果可得：
% 这种对 $\phi_\mathrm{a}$ 的在相区宽度增幅比中同样存在。
% 由图~\ref{fig:PhaseWidth}a 可见，当以 $\phi_\mathrm{a}=0.25$ 为基准进行比较时，$\phi_\mathrm{a}=0.50$ 的结果表明：
% \begin{equation}
%     \begin{split}
%         \frac{\Delta \bar{\phi}_{\mathrm{A}}(\phi_{\mathrm{a}}=0.50,\mathrm{C})}
%         {\Delta \bar{\phi}_{\mathrm{A}}(\phi_{\mathrm{a}}=0.25,\mathrm{C})} <&
%         \frac{\Delta \bar{\phi}_{\mathrm{A}}(\phi_{\mathrm{a}}=0.50,\mathrm{BCC})}
%         {\Delta \bar{\phi}_{\mathrm{A}}(\phi_{\mathrm{a}}=0.25,\mathrm{BCC})} \\
%         <&\frac{\Delta \bar{\phi}_{\mathrm{A}}(\phi_{\mathrm{a}}=0.50,\mathrm{FCC})}
%         {\Delta \bar{\phi}_{\mathrm{A}}(\phi_{\mathrm{a}}=0.25,\mathrm{FCC})} 
%     \end{split}
% \end{equation}
% 该不等式表明，滑动自由度对界面复杂度较高相的稳定化作用更强，且其对 $\phi_{\mathrm{a}}$ 变化的响应也更为显著；特别是在 $\phi_\mathrm{a}\approx 0.50$ 附近，这种作用在 FCC 相中最为突出，表明在 $N\chi$ 较高时，L-S TT 有利于球状相的形成。
By comparing the variation of phase widths at different $N\chi$, we can clarify the competition between entropic and enthalpic effects in the L--S TT system.
Fig.~\ref{fig:PhaseWidth}a shows how the widths of the ordered phase regions change with increasing $N\chi$ for different $\phi_{\mathrm a}$.
The stacked bar plot gives the relative contribution of each phase to the total ordered-region width, whereas the upper line plot compares the normalized width $\Delta\bar{\phi}_{\mathrm A}(\phi_{\mathrm a},X)$ for different $\phi_{\mathrm a}$.
Here, $X\in\{\mathrm{L},\mathrm{G},\mathrm{C},\mathrm{BCC},\mathrm{FCC}\}$ denotes the ordered phase.
Overall, the total width of the ordered region increases with $N\chi$, consistent with the general tendency of block copolymers to form ordered structures when the incompatibility between blocks becomes stronger.
However, this increase depends strongly on $\phi_{\mathrm a}$.

We interpret this dependence through the coupling between the conformational response to molecular architecture and the interfacial requirements of the ordered phases.
The effective conformational-entropy scale in Fig.~\ref{fig:ODT}e varies nonmonotonically with $\phi_{\mathrm a}$, providing a reference for this interpretation rather than a direct entropy decomposition for each ordered phase.

The response to $\phi_{\mathrm a}$ also depends on the interfacial geometry of each ordered phase.
The L phase has regular flat interfaces, but it requires nearly balanced A and B volume fractions, a condition not generally satisfied by L--S TT copolymers; therefore, conformational regulation provides only a limited stabilization effect, leading to a nonmonotonic change in its phase fraction, as shown in Fig.~\ref{fig:PhaseWidth}a.
The G phase contains interpenetrating network channels and is more sensitive to flexible chain regulation, so its phase fraction decreases markedly as $\phi_{\mathrm a}$ increases.
By contrast, the C phase and the spherical phases, including BCC and FCC, have more complex and highly curved interfaces and more constrained packing environments.
They therefore rely more strongly on the ability of the sliding degree of freedom to regulate the density field, and their phase widths reach maxima near $\phi_{\mathrm a}\approx0.50$.
This phase-dependent response is also reflected in the relative increase of phase widths.
Using $\phi_{\mathrm a}=0.25$ as the reference, the results at $\phi_{\mathrm a}=0.50$ in Fig.~\ref{fig:PhaseWidth}a show that
\begin{equation}
    \begin{split}
        \frac{\Delta \bar{\phi}_{\mathrm{A}}(\phi_{\mathrm{a}}=0.50,\mathrm{C})}
        {\Delta \bar{\phi}_{\mathrm{A}}(\phi_{\mathrm{a}}=0.25,\mathrm{C})} <
        \frac{\Delta \bar{\phi}_{\mathrm{A}}(\phi_{\mathrm{a}}=0.50,\mathrm{BCC})}
        {\Delta \bar{\phi}_{\mathrm{A}}(\phi_{\mathrm{a}}=0.25,\mathrm{BCC})} 
        <
        \frac{\Delta \bar{\phi}_{\mathrm{A}}(\phi_{\mathrm{a}}=0.50,\mathrm{FCC})}
        {\Delta \bar{\phi}_{\mathrm{A}}(\phi_{\mathrm{a}}=0.25,\mathrm{FCC})}.
    \end{split}
\end{equation}
This inequality indicates that the stabilizing effect of the sliding degree of freedom becomes stronger for phases with more complex interfacial geometry.
The effect is strongest for the FCC phase, suggesting that at relatively high $N\chi$, L--S TT copolymers favor the formation of spherical phases.


% 为了阐明不同拓扑调节机制在复杂有序结构形成中的作用，我们考察了图~\ref{fig:PhaseWidth}b,e 所示的 BCC 与 FCC 相区宽度 $\Delta \bar{\phi}_{\mathrm{A}}$ 随 $N\chi$ 和 $\phi_\mathrm{a}$ 的演化行为。
% 结果表明，在所有 $\phi_\mathrm{a}$ 下，BCC 相区宽度均随 $N\chi$ 单调增加，这说明随着焓的增加，BCC 相始终能够稳定受益于滑动自由度$f$的引入所带来的构象调节。
% 结合不同 $\phi_\mathrm{a}$ 条件下各等值面 $\phi_{\mathrm A}^{\mathrm{iso}}$ 上的 $f_\Omega^\ast$ 的分布特征可解释该现象。
% 如图~\ref{fig:PhaseWidth}c，当 $\phi_\mathrm{a}=0.50$ 时，$f_\Omega^\ast$ 随 $\phi_{\mathrm A}^{\mathrm{iso}}$ 的变化趋势在不同 $N\chi$ 下表现出明显差异。
% 在较低 $N\chi$ 条件下，$f_\Omega^\ast$ 在 B 组分胶束中心趋近于 $0$，表明 L-S TT 更倾向于线型 diblock 构型；而在较高 $N\chi$ 条件下，胶束中心处的 $f_\Omega^\ast$ 则保持为非零的有限值，说明该区域出现了对星型 A$_2$AB 构型的偏好。
% 此外，在 A 组分作为 shell 的组成区域中，$f_\Omega^\ast$ 的取值也随 $N\chi$ 的增大而逐渐降低。
% 图~\ref{fig:PhaseWidth}d 进一步表明，在固定 $\phi_{\mathrm A}^{\mathrm{iso}}=0.75$ 时，最高 $N\chi$ 对应的 $f_\Omega^\ast$ 曲线近似为一条非零水平线，说明在强分离条件下，该区域的最优拓扑选择已基本摆脱对 $\phi_{\mathrm{a}}$ 的依赖，并锁定在一个稳定的非零值附近。
% 这些结果表明，在低 $N\chi$ 区域，体系的拓扑选择主要受滑动熵调控。
% 随着 $N\chi$ 增加，界面相关的焓效应逐步增强，并不断削弱这种由熵主导的拓扑择优，使 $f$ 的分布由以整体构象熵调节为主，逐渐转向以缓解局部 packing frustration 为主，进而表现为对 $\phi_\mathrm{a}$ 依赖的减弱，以及在不同空间位置（$\phi_{\mathrm A}^{\mathrm{iso}}$）上的功能分化增强。
% 需要指出的是，在本研究考察的 $N\chi\le 45.0$ 范围内，这种焓主导趋势尚未完全压制滑动自由度所带来的构象调节收益，因此 BCC 相仍能够持续从拓扑调节中获益，并保持稳定扩张的相区宽度。
To clarify how different topology-regulation mechanisms contribute to the formation of complex ordered structures, we examine the evolution of the BCC and FCC phase widths, $\Delta\bar{\phi}_{\mathrm A}$, with $N\chi$ and $\phi_{\mathrm a}$, as shown in Fig.~\ref{fig:PhaseWidth}b,e.
For all $\phi_{\mathrm a}$, the BCC phase width increases monotonically with $N\chi$, indicating that BCC consistently benefits from the conformational regulation enabled by the sliding degree of freedom.
This behavior can be understood from the distribution of the preferred topology $f_\Omega^\ast$ on different density layers $\phi_{\mathrm A}^{\mathrm{iso}}$.
As shown in Fig.~\ref{fig:PhaseWidth}c, for $\phi_{\mathrm a}=0.50$, the variation of $f_\Omega^\ast$ with $\phi_{\mathrm A}^{\mathrm{iso}}$ depends strongly on $N\chi$.
At lower $N\chi$, $f_\Omega^\ast$ approaches zero near the B-rich micellar center, indicating a preference for an AB linear diblock configuration.
At higher $N\chi$, however, $f_\Omega^\ast$ remains finite in the same region, showing that an A$_2$AB-like topology becomes locally favored.
In the A-shell-forming composition range, $f_\Omega^\ast$ also decreases gradually as $N\chi$ increases.
Fig.~\ref{fig:PhaseWidth}d further shows that, at fixed $\phi_{\mathrm A}^{\mathrm{iso}}=0.75$, the $f_\Omega^\ast$ curve at the highest $N\chi$ is nearly horizontal and remains at a finite value.
This indicates that, under strong segregation, the preferred topology in this region becomes largely independent of $\phi_{\mathrm a}$ and is locked around a stable nonzero value.
These results suggest a crossover in the role of sliding topology.

At low $N\chi$, the topology preference remains close to the RPA--MUT reference, reflecting the stronger role of the conformational response associated with the sliding ensemble.
As $N\chi$ increases, stronger interfacial and packing constraints weaken this global topology preference and promote increasingly local topology regulation in response to heterogeneous filling environments, as reflected by the reduced dependence on $\phi_{\mathrm a}$ and the stronger functional differentiation across spatial layers.

Within the range considered here, $N\chi\le45.0$, the stronger interfacial constraints do not eliminate the benefit of local topology adjustment.
Consequently, the BCC phase continues to gain stability from topology regulation and its phase width expands steadily.


% 相比 BCC 相，图~\ref{fig:PhaseWidth}e 显示 FCC 相区宽度 $\Delta \bar{\phi}_{\mathrm{A}}$ 随 $N\chi$ 呈现明显的非单调变化：在中等分离强度 $N\chi \approx 35$ 附近达到最大，而在更高 $N\chi$ 下则略有收缩。
% 这表明，滑动自由度对 FCC 相的稳定化作用只在有限的分离强度范围内最为有效；当界面相关焓贡献进一步增强时，这种拓扑调节收益反而开始减弱。
% 结合图~\ref{fig:PhaseWidth}f,g 中 $f_\Omega^\ast$ 的分布特征可以理解这一行为。
% 如图~\ref{fig:PhaseWidth}f 所示，在 $\phi_\mathrm{a}=0.50$ 时，FCC 相中 $f_\Omega^\ast$ 对空间位置 $\phi_{\mathrm A}^{\mathrm{iso}}$ 的响应主要集中在 A shell 对应的高 $\phi_{\mathrm A}^{\mathrm{iso}}$ 区域，而在其余区域内始终为 0，即维持 AB 线型 diblock 构型。
% 这说明，相较于 BCC 相中较为连续的空间调节，FCC 相中的滑动拓扑调节已被压缩到更狭窄的局域区域，其可参与构象优化的空间范围明显受限。
% 图~\ref{fig:PhaseWidth}g 表明，在固定 $\phi_{\mathrm A}^{\mathrm{iso}}=0.75$ 时，$f_\Omega^\ast$ 随 $\phi_\mathrm{a}$ 的变化同样呈现出与 BCC 相相似的 $N\chi$ 依赖行为。
% 因此，FCC 相区所表现出的极值行为可理解为熵效应与焓效应短暂平衡的结果。
% 在中等 $N\chi$ 下，滑动自由度仍能在界面附近提供有限但有效的拓扑重排，从而帮助体系适应 FCC 所要求的更复杂局域填充与曲率约束，因此 FCC 相区得以扩张。
% 而当 $N\chi$ 进一步增大时，界面约束和局域 packing frustration 显著增强，导致体系已难以通过进一步的拓扑分化来持续降低界面自由能，最终表现为 FCC 相稳定性下降及其相区宽度的回缩。
% 总体而言，BCC 与 FCC 的差异表明，滑动自由度对有序相的稳定化作用并非简单随界面复杂度单调增强，而是取决于该结构是否允许拓扑调节在足够大的空间范围内持续分担局域填充与界面匹配任务。
% 这一现象完整呈现了 L-S TT 体系中熵和焓两种机制之间的竞争和协同关系。
Compared with BCC, the FCC phase shows a nonmonotonic dependence of the composition width $\Delta\bar{\phi}_{\mathrm A}$ on $N\chi$ (Fig.~\ref{fig:PhaseWidth}e).
The FCC composition window reaches its maximum width near $N\chi\approx35$ and then slightly narrows at higher $N\chi$, while FCC remains stable at selected compositions up to $N\chi=45$ within the parameter range examined here. 


The contrasting BCC and FCC phase widths are accompanied by different patterns of topology preference across the density layers.
For soft polymer domains, the more nearly spherical Voronoi cell of BCC can reduce the geometric distortion required for nominally spherical domains to fill space.
Sliding provides an additional means of adapting chain conformations to the local filling requirements of each structure.


Fig.~\ref{fig:PhaseWidth}f,g resolves this topology response in FCC.
For $\phi_{\mathrm a}=0.50$, the response of $f_\Omega^\ast$ to the spatial layer $\phi_{\mathrm A}^{\mathrm{iso}}$ in FCC is mainly confined to the high-$\phi_{\mathrm A}^{\mathrm{iso}}$ region associated with the A shell, whereas $f_\Omega^\ast$ remains close to zero in the other regions, corresponding to the AB linear diblock configuration.
Compared with the more continuous variation of the preferred topology in BCC, this response in FCC is concentrated within a narrower range of density layers.
Fig.~\ref{fig:PhaseWidth}g further shows that, at fixed $\phi_{\mathrm A}^{\mathrm{iso}}=0.75$, the dependence of $f_\Omega^\ast$ on $\phi_{\mathrm a}$ follows an $N\chi$-dependent trend similar to that observed in BCC.
Together, these trends suggest that the maximum in the FCC phase width reflects a competition between local conformational adaptation and stronger interfacial constraints as $N\chi$ increases.
The entropy loss in Fig.~\ref{fig:SpherePOmega}c provides a quantitative comparison of the accompanying topology homogenization.
At $N\chi=45$ and $\phi_{\mathrm a}=0.50$, the representative A-shell states exhibit a larger sliding-distribution entropy loss in BCC than in FCC.
The smaller value for FCC indicates a weaker departure of the local topology distributions from uniform sliding under the A segment weighting, and hence stronger topology homogenization in this representative state.
The same BCC--FCC ordering is also observed for the representative A core states.
Redistribution of topology weights may help accommodate the local filling requirements of FCC at intermediate $N\chi$, whereas its subsequent narrowing indicates that competing phases become favorable over a larger composition range at higher $N\chi$.
The contrast with the steadily expanding BCC window suggests that the stabilizing effect of sliding depends on how effectively the accessible topology response adapts to the distinct local filling geometry of each structure.


